\begin{table}[t]
  \centering
  \caption{\textbf{古典的コートモデル当てはめの補助 stress probe。}
  \cref{tab:court-alignment} とは別の決定論的な標本
  （\resClassicProbeSamples）で固定版を走らせた結果であり、
  主比較の数値と直接比較しない。
  実装は \cite{gchlebusrepo} の commit \resClassicCommitShort を
  OpenCV 4 向けに patch（\texttt{12422aec}）してビルドし、
  既定閾値のまま CPU のみで実行した。
  成功とは有限な 16 点を出力したことを指し、
  幾何の正しさは含まない。
  }
  \label{tab:classical_probe}
  \small
  \begin{tabularx}{\linewidth}{@{}p{0.20\linewidth}r r X@{}}
    \toprule
    標本 & 実行 & 出力 16 点が有限 & 観察 \\
    \midrule
    実写 10 件 &
      \resClassicRealSuccess & 全件 &
      \resClassicRealInFramePartial 件は \resClassicRealInFramePairs 点のみ画像内で、
      大きな外挿を含む \\
    合成 \texttt{B00} 10 件 &
      \resClassicSynthSuccess & 成功分のみ &
      失敗 \resClassicSynthFailure（クラッシュ \resClassicSynthCrash，
      exit 3 が \resClassicSynthExitThree）。
      成功分にも誤フィットと縮退がある \\
    \bottomrule
  \end{tabularx}
\end{table}
